\subsection{Steered OPD (SOPD) from textual feedback}

Composer 2.5 calls its method \emph{targeted textual feedback}; technically it is
privileged-context on-policy self-distillation, often abbreviated OPSD. Here, \emph{steered
OPD} (SOPD) refers to this specific use of a hint, not Cursor's name for the method.
Self-Distillation Fine-Tuning uses the same general construction with a
demonstration-conditioned copy of the student as its teacher
\citep{shenfeld2026selfdistillation}.
Suppose one turn in a long coding rollout calls an unavailable tool. Its effect may
disappear inside the terminal reward. At that turn, a grounded hint---derived from an
environment error or a generated diagnosis---is inserted only into the teacher context.
The original-context policy remains the student, and a KL loss moves the student's
distribution toward the hint-conditioned distribution for that turn. The broader
trajectory still receives its RL objective \citep{cursor2026composer25}. The hint therefore
supplies a semantic correction without becoming an inference-time input.

A proposed extension is to internalize the feedback process itself. Instead of only
distilling the hidden hint's effect, train the model on sequences of
\emph{state $\rightarrow$ grounded hint $\rightarrow$ revised action}, so it learns to
produce and use a diagnosis at inference time. SFT on hint text alone is insufficient: the
model could learn to verbalize a good reminder without changing its next decision. The
hint must be coupled to a verified correction, and self-generated hints must be checked to
avoid reinforcing plausible but false diagnoses. This explicit route costs inference
tokens but may create reusable trace-level skills; steered OPD is cheaper at inference
because it compiles the hint's behavioral effect directly into the policy.
